\documentclass[11pt]{article}

\usepackage[a4paper, margin=2.5cm]{geometry}
\usepackage{amsmath}
\usepackage{amssymb}
\usepackage{booktabs}
\usepackage{graphicx}
\usepackage{parskip}
\usepackage[hidelinks]{hyperref}

\title{A Focus-Cost Model for Meeting Placement}
\author{}
\date{}

\begin{document}
\maketitle

\section{Motivation}

Meeting fatigue is usually attributed to the number of meetings. This
document sets out a model in which the dominant factor is instead
\emph{fragmentation}: where meetings sit within a day, rather than how much
of the day they occupy.

Two hours of meetings placed at 09:00 and 15:00 leave nobody with more than
three hours of uninterrupted time. The same two hours booked 09:00 to 11:00
leave a five hour block. The load is identical; the days are not.

\section{Notation}

All quantities are defined for a single person on a single day. The day is
the unit of computation because the interactions that matter, namely
fragmentation, fatigue and the loss of a lunch break, occur within a day and
not across days.

\subsection*{Sets}

\begin{tabular}{@{}ll@{}}
\toprule
$W = [w_s, w_e]$ & working window \\
$M$ & the person's meetings on the day \\
$G(M)$ & gaps: free intervals within $W$ induced by $M$ \\
$R(M)$ & runs: maximal contiguous stretches of meeting time \\
$A(m)$ & attendee list of a candidate meeting $m$ \\
$N$, $n = |N|$ & meetings in a person-day, used for attribution \\
\bottomrule
\end{tabular}

\medskip

Gaps and runs are derived from $M$, not stored. Overlapping meetings merge
into a single run.

\subsection*{Parameters}

\begin{tabular}{@{}llll@{}}
\toprule
Symbol & Meaning & Default & Kind \\
\midrule
$c$ & re-entry cost & 20 min & measurement \\
$\lambda$ & fatigue rate & 0.4 & measurement \\
$r^{*}$ & fatigue threshold & 90 min & measurement \\
$\mu$ & lunch penalty & 60 min & policy \\
$L$ & lunch window & 12:00--14:00 & policy \\
$\ell$ & required lunch length & 60 min & policy \\
\bottomrule
\end{tabular}

\medskip

Measurement parameters are intended to be fitted against self-reported data
and are not asserted. Policy parameters encode a norm the organisation
chooses and are not subject to calibration. $\lambda$ is dimensionless; all
other parameters are in minutes.

\section{The model}

\subsection{Gap yield}

A free interval is not worth its face value, since some minutes at the start
of any free stretch are lost re-entering the task:
\[
u(g) \;=\; \max\bigl(0,\; |g| - c\bigr)
\]

A 60-minute gap yields 40 minutes of usable focus. A 15-minute gap yields
nothing at all. The re-entry cost is paid once per gap, so a day broken into
four pieces pays it four times.

\subsection{Fatigue}

Long unbroken runs of meetings degrade what follows them:
\[
\Phi(M) \;=\; \lambda \sum_{r \in R(M)} \max\bigl(0,\; |r| - r^{*}\bigr)
\]

This is convex in run length, so the fourth consecutive meeting is charged
more than the first.

\subsection{Lunch protection}

A day leaving no adequate break within the lunch window is penalised. The
check is against raw gap lengths rather than gap yields, and is all or
nothing:
\[
\Lambda(M) \;=\; \mu \cdot \mathbf{1}\Bigl[\, \nexists\, g \in G(M) \;:\;
g \subseteq L \;\wedge\; |g| \ge \ell \,\Bigr]
\]

\subsection{Usable focus time and the day penalty}

Write $Y(M)$ for the total yield of the day's free intervals:
\[
Y(M) \;=\; \sum_{g \in G(M)} \max\bigl(0,\; |g| - c\bigr)
\]

Usable focus time subtracts the two charges from that yield:
\[
U(M) \;=\; Y(M) \;-\; \Phi(M) \;-\; \Lambda(M)
\]

An empty day contains a single gap spanning $W$, so
\[
U(\varnothing) \;=\; |W| - c
\]
and the day penalty is the shortfall against it:
\[
P(M) \;=\; U(\varnothing) - U(M)
\]

$P$ is denominated in minutes of focus time destroyed. An empty day scores
zero; a fully booked day scores $|W| - c$.

\subsection{Collected form}

Expanding every term gives the penalty function in full, as three
contributions to the shortfall against a clear day:

\begin{center}
\resizebox{\textwidth}{!}{$\displaystyle
P(M) \;=\; \bigl(|W| - c\bigr)
\;+\; \underbrace{\sum_{g \in G(M)} -\max\bigl(0,\; |g| - c\bigr)}_{\text{fragmentation}}
\;+\; \underbrace{\lambda \sum_{r \in R(M)} \max\bigl(0,\; |r| - r^{*}\bigr)}_{\text{fatigue}}
\;+\; \underbrace{\mu \cdot \mathbf{1}\Bigl[\, \nexists\, g \in G(M) : g \subseteq L \;\wedge\; |g| \ge \ell \,\Bigr]}_{\text{lunch}}
$}
\end{center}

The fragmentation term is the negative of the day's focus yield. The more
usable focus the day retains, the smaller the shortfall, so a day broken
into more pieces contributes more to $P$ through paying $c$ once per gap.

\section{Marginal cost}

The cost of a candidate meeting is the increase in day penalty it causes,
summed over everyone invited:
\[
\mathrm{Cost}_p(m \mid M) \;=\; P\bigl(M \cup \{m\}\bigr) - P(M)
\qquad
\mathrm{Cost}(m) \;=\; \sum_{p \in A(m)} \mathrm{Cost}_p\bigl(m \mid M_p\bigr)
\]

Since $U(\varnothing)$ cancels between the two evaluations, this decomposes
into the change in each component:
\[
\mathrm{Cost}_p(m \mid M) \;=\;
\underbrace{\bigl[\, Y(M) - Y(M \cup \{m\}) \,\bigr]}_{\text{focus yield lost}}
\;+\; \underbrace{\bigl[\, \Phi(M \cup \{m\}) - \Phi(M) \,\bigr]}_{\text{fatigue added}}
\;+\; \underbrace{\bigl[\, \Lambda(M \cup \{m\}) - \Lambda(M) \,\bigr]}_{\text{lunch added}}
\]

Marginal cost is a function of calendar state, not a property of the
meeting. It is not additive: subadditive when meetings abut, since the
second pays only its duration, and superadditive once a combined run exceeds
$r^{*}$.

\section{Closed form under the fragmentation term}

The results in this section hold when $\Phi$ and $\Lambda$ do not bind. Both
additional terms deliberately break them.

For a candidate of duration $D$ placed inside a single gap of length $g$,
leaving fragments $a$ on the left and $b$ on the right with $a + D + b = g$:
\[
\mathrm{Cost} \;=\; \max(0, g-c) \;-\; \max(0, a-c) \;-\; \max(0, b-c)
\;=\;
\begin{cases}
D + c, & a \ge c \;\wedge\; b \ge c \\[2pt]
D + a, & a < c \;\wedge\; b \ge c \\[2pt]
D + b, & a \ge c \;\wedge\; b < c \\[2pt]
g - c, & a < c \;\wedge\; b < c
\end{cases}
\]

Cost is therefore constant at $D + c$ across the interior of a gap and falls
linearly to $D$ at its edges. The fourth case gives a cost below the
meeting's own duration, which is correct: a candidate filling a gap entirely
consumes time that was partly lost to re-entry regardless.

\subsection{Spread}

It follows that when a flush position exists for every attendee, the
difference between the dearest and cheapest bookable slot is
\[
\max_{s} \mathrm{Cost}(m_s) \;-\; \min_{s} \mathrm{Cost}(m_s)
\;=\; |A| \cdot c
\]
independent of duration and of the shape of the day. Each attendee pays
either $D$ or $D + c$, so the difference is $c$ per attendee.

This is a prediction rather than an observation, and it was confirmed on 13
of 20 randomly generated scenarios. The remaining seven lacked a candidate
position leaving every attendee flush, and so never reached the lower bound.

\section{Retrospective attribution}

Marginal cost depends on the order in which meetings were booked, which is
invisible on a calendar, so it cannot be used to rank meetings against one
another. For a person-day with meeting set $N$, allocate the day penalty by
Shapley value: each meeting's marginal contribution, averaged over every
order in which the meetings could have arrived.
\[
\varphi_i \;=\; \sum_{S \subseteq N \setminus \{i\}}
\frac{|S|!\,\bigl(n - |S| - 1\bigr)!}{n!}
\Bigl[\, P\bigl(S \cup \{i\}\bigr) - P(S) \,\Bigr]
\]

This is the unique allocation satisfying efficiency, symmetry, the null
player property and additivity. Efficiency,
\[
\sum_{i \in N} \varphi_i \;=\; P(N),
\]
is the property that matters in practice: attributions across a day sum
exactly to that day's penalty, so per-meeting figures reconcile with
per-person totals.

Exact computation requires $n \cdot 2^{\,n-1}$ evaluations of $P$, which is
negligible for the meeting counts found in a single day. Above roughly
$n = 12$ the implementation samples permutations by Monte Carlo and reports
a standard error.

\subsection{Which measure belongs where}

Marginal cost is prospective and answers \emph{what will this booking cost},
so it belongs in the scheduling view, where state-dependence is information
rather than distortion. Shapley attribution is retrospective and answers
\emph{what share of this day's damage does this meeting own}, so it belongs
on the calendar, where the figures are expected to sum.

\section{Limitations}

\textbf{Linear gap yield.} $u(g)$ treats a minute of free time as equally
valuable regardless of the block it sits in, so breaking a four-hour block
and breaking a forty-minute gap are priced similarly. Weighting yield by
block quality, for instance $u(g) = (g - c)\min(1, g/g^{*})$, would
distinguish them at the cost of one further parameter and of the closed form
above.

\textbf{No benefit term.} The model prices cost and never value. Its
unconstrained optimum is to hold no meetings, which is plainly wrong. A
sampled post-meeting value signal would give a cost-versus-value view, but
requires participation in a way the cost side does not.

\textbf{Uncalibrated parameters.} The measurement parameters are starting
values. The intended fitting procedure is a grid search maximising Spearman
rank correlation between $P$ and participants' own rankings of their past
weeks.

\end{document}
